\scanpage{181}
\mspage{55}
% archive sheet number: 180

(5.4) n'exclut nullement que $X_0$ admette un projecteur $p$ ($p^2=p$) non trivial. Celui-ci aura une image $X'$ \deleted{dans} $\mathcal M$, qui sera facteur direct de $X_0$ non isom. à $X_0$ i.e. non dans $\deleted{\unsure{M}}M$.

Ces difficultés disparaissent si on suppose que $M^{\unsure{\pi|\Pi}}$ est stable par $\unsure{\underline{\lim}|\varinjlim}$ filtrantes indexées par des ens. \deleted{\unsure{ord}} ordonnés $I$ de cardinal $\leq\pi$. En effet, dans ce cas, \deleted{\unsure{le sous}} dans la démonstration précédente, le sous-ensemble $I_M$ est stable par bornes supérieures (dans $I$) des parties $J$ de $I_M$ filtrantes et de cardinal $\leq\pi$. Donc le syst. inductif $(\unsure{X|X'}_i)_{i\in I_M}$ est encore $\pi$-adapté, et on peut lui appliquer le Lemme 5.7.) pour trouver un sous-ensemble cofinal $I'_M\subset I_M$, tel que les $X_i\unsure{\longrightarrow|-} X$ pour $i\in I'_M$ soient monos. L'image $J$ de $I'_M$ dans $\mathrm{Ss}_{\pi}(X)$ est alors une partie de $\mathrm{Ss}_{\pi}(X)$ telle que $\varinjlim_{\alpha\in J} Z_\alpha=X$, donc (par l'hyp. faite sur $\pi$) cofinale dans $\mathrm{Ss}_{\pi}(X)$ (5.4.), avec $Z_\alpha\in M^\pi$ \deleted{\unsure{qq}} pour tt $\alpha\in J$. Donc par (5.3.) on a $X\in M$, qed.

